% SPDX-FileCopyrightText: © 2026 Idiap Research Institute <contact@idiap.ch>
%
% SPDX-FileContributor: Andre Freitas <andre.freitas@idiap.ch>
% Neuro-symbolic AI Group
%
% SPDX-FileContributor: Danilo Gusicuma <danilo.gusicuma@idiap.ch>
% Neuro-symbolic AI Group
%
% SPDX-License-Identifier: MIT
\documentclass[11pt]{article}

\usepackage[margin=1in]{geometry}

\usepackage{amsmath,amssymb,amsthm}

\title{Dormant Paper-Argument Draft for Yang--Mills Milestone Certification}

\author{Abducer Long-Horizon Monitor}

\date{}

\newtheorem{claim}{Claim}

\begin{document}

\maketitle

\begin{abstract}

No monitored milestone has reached the archive trigger. This document records the argument shape that will be filled only after certification evidence is attached.

\end{abstract}

\section{Claim Scope}

A paper argument is activated only for milestones with certification evidence attached to the same continuum Yang-Mills object: the gauge-field construction, OS reconstruction, nontriviality witness, and spectral-gap witness must all be read from one limiting theory. Until then, the draft records the intended argument shape without claiming a proof.

\section{Argument Spine}

\begin{itemize}
\item State the reached milestone and its exact certification evidence.
\item Identify the same-theory gauge-field object carrying construction, reconstruction, nontriviality, and spectral content.
\item Present the mathematical bridge lemmas used to move from Euclidean Schwinger data and gauge-invariant observables to the milestone claim.
\item List disqualifier tests and how each was ruled out.
\item State remaining obligations and the boundary between achieved progress and unresolved Clay target commitments.
\end{itemize}

\section{Activation Condition}

The draft remains dormant until at least one milestone has status achieved, potentially\_achieved, potential\_breakthrough\_signal, or certified\_breakthrough.

\section{Current Non-Claim}

The current Yang-Mills run has useful bridge-lemma packets and monitoring artifacts, but no big-deal milestone is certified as reached.

\section{Physics Mechanism Learning To Preserve}

\begin{itemize}
\item We learned that Hamiltonian language is indispensable as an endpoint test but weak as a discovery engine unless the reconstruction inputs are already available.
\item A claimed gap must always be asked backwards: where did H\_YM, Omega, H, P\_perp, and the observable algebra come from?
\item This converted the analysis from naming a desired spectral inequality into tracking the Euclidean construction and reconstruction prerequisites that make the inequality meaningful.
\item Correlator decay is valuable only when it is a same-theory spectral witness: the observables, OS reconstruction, and Hamiltonian must all come from the same limiting Schwinger functions.
\item The right mathematical object to demand next is not just a decay estimate, but a cutoff-uniform decay-to-spectrum lemma with an observable-density hypothesis.
\item This taught the loop to treat strong-coupling lattice decay as evidence-task material until it carries a continuum-transfer certificate.
\item Balaban-style control should be used as the construction spine, not as a substitute for the mass-gap mechanism.
\item The decisive test is whether RG-controlled cutoff objects produce OS-positive continuum Schwinger functions and a non-Gaussian same-theory witness.
\item This shifted the loop toward an RG-to-OS bridge lemma instead of adding more renormalisation vocabulary.
\item The analysis should now advance by testing typed bridge lemmas, not by accumulating more route families.
\item A strategy only counts as learning if it changes the mathematical dependency order: Hamiltonian claims move behind reconstruction gates, decay claims move behind uniformity/density gates, and RG claims move behind OS/nontriviality gates.
\item The next useful artifact is therefore a bridge-lemma packet with definitions, assumptions, equations, failure tests, and conditional skill triggers.
\item The loop learned to use contradiction as a disciplined pressure test rather than an excuse to skip construction.
\item A branch succeeds only if the contradiction target is accepted or validator-backed in the same theory.
\item A failed branch is still valuable when it tells us whether to pivot toward density, closedness, nontriviality, or uniform-decay evidence.
\end{itemize}

\section{Gedanken Experiment Boundary}

Gedanken packets may affect branch control only through validated follow-up queue items, transfer verdicts, and qualitative graph annotations. They are not counted as proof progress unless a later derivation, computation, or source-backed certificate discharges the relevant obligation.

\section{Proof-Safety Boundary}

The paper draft must distinguish certified progress from a full Yang-Mills proof unless all Clay commitments are discharged in the same theory.

\end{document}
